\documentclass[12pt,a4paper]{exam}
\usepackage[utf8]{inputenc}
\usepackage{amsmath,amssymb,amsfonts}
\usepackage{geometry}
\usepackage{enumitem}
\usepackage{graphicx}
\usepackage{hyperref}
\usepackage{xcolor}
\geometry{a4paper, top=2cm, bottom=2cm, left=2.5cm, right=2.5cm}
\hypersetup{colorlinks=true, linkcolor=blue, urlcolor=blue}
\renewcommand{\thequestion}{\textbf{\arabic{question}}}
\renewcommand{\questionlabel}{\thequestion.}
\pointname{ marks}
\pointsdroppedatright
\bracketedpoints
\setlength{\parindent}{0pt}
\setlength{\emergencystretch}{3em}
\allowdisplaybreaks
\newsavebox{\schmathbox}
\newcommand{\fitmath}[1]{%
  \sbox{\schmathbox}{$\displaystyle #1$}%
  \ifdim\wd\schmathbox>\linewidth
    \resizebox{\linewidth}{!}{\usebox{\schmathbox}}%
  \else
    \usebox{\schmathbox}%
  \fi}

\begin{document}

\thispagestyle{empty}
\begin{center}
  \textbf{\LARGE Diag} \\[0.3cm]
  \textbf{Time Allowed: 3 Hours} \hfill \textbf{Maximum Marks: 80} \\[0.3cm]
  \rule{\textwidth}{0.5pt}
\end{center}

\vspace{0.3cm}

\noindent\textbf{General Instructions:}
\begin{enumerate}
  \item {\normalfont\normalcolor This question paper contains 40 questions. All questions are compulsory.}
  \item {\normalfont\normalcolor Questions are grouped into sections A through E.}
  \item {\normalfont\normalcolor Use of calculators is NOT permitted.}
\end{enumerate}
\bigskip
\noindent\textbf{\large Section A: Multiple Choice \& Assertion-Reason (1 mark each)}

\begin{questions}
\question[1] To draw a pair of tangents to a circle which are inclined to each other at an angle of $60^{\circ}$, it is required to draw tangents at the endpoints of those two radii of the circle, the angle between which is:
\begin{enumerate}[label=(\Alph*)]
  \item (A) $60^{\circ}$
  \item (B) $90^{\circ}$
  \item (C) $120^{\circ}$
  \item (D) $180^{\circ}$
\end{enumerate}
\vspace{0.15cm}

\question[1] If a point $P$ is $10$ cm away from the centre of a circle of radius $6$ cm, the length of the tangent drawn from $P$ to the circle is:
\begin{enumerate}[label=(\Alph*)]
  \item (A) $7$ cm
  \item (B) $8$ cm
  \item (C) $9$ cm
  \item (D) $10$ cm
\end{enumerate}
\vspace{0.15cm}

\question[1] How many tangents can be drawn from an external point to a circle?
\begin{enumerate}[label=(\Alph*)]
  \item (A) 1
  \item (B) 2
  \item (C) 0
  \item (D) Infinite
\end{enumerate}
\vspace{0.15cm}

\question[1] To divide a line segment $AB$ in the ratio $3:4$, we draw a ray $AX$ such that $\angle BAX$ is acute. Then we mark points on $AX$ at equal distances. The minimum number of points to be marked is:
\begin{enumerate}[label=(\Alph*)]
  \item (A) 3
  \item (B) 4
  \item (C) 7
  \item (D) 12
\end{enumerate}
\vspace{0.15cm}

\question[1] The angle between a tangent to a circle and the radius passing through the point of contact is:
\begin{enumerate}[label=(\Alph*)]
  \item (A) $45^{\circ}$
  \item (B) $60^{\circ}$
  \item (C) $90^{\circ}$
  \item (D) $180^{\circ}$
\end{enumerate}
\vspace{0.15cm}

\question[1] If the radius of a sphere is $7 \text{ cm}$, then its surface area is:
\begin{enumerate}[label=(\Alph*)]
  \item (A) 308 cm\textasciicircum{}2
  \item (B) 440 cm\textasciicircum{}2
  \item (C) 616 cm\textasciicircum{}2
  \item (D) 1232 cm\textasciicircum{}2
\end{enumerate}
\vspace{0.15cm}

\question[1] The volume of a right circular cylinder with radius $3 \text{ cm}$ and height $7 \text{ cm}$ is:
\begin{enumerate}[label=(\Alph*)]
  \item (A) 154 cm\textasciicircum{}3
  \item (B) 198 cm\textasciicircum{}3
  \item (C) 264 cm\textasciicircum{}3
  \item (D) 396 cm\textasciicircum{}3
\end{enumerate}
\vspace{0.15cm}

\question[1] Assertion (A): The length of the tangent drawn from an external point to a circle is always equal to the radius of the circle. Reason (R): The tangent at any point of a circle is perpendicular to the radius through the point of contact.
\begin{enumerate}[label=(\Alph*)]
  \item (A) Both Assertion (A) and Reason (R) are true and Reason (R) is the correct explanation of Assertion (A)
  \item (B) Both Assertion (A) and Reason (R) are true but Reason (R) is NOT the correct explanation of Assertion (A)
  \item (C) Assertion (A) is true but Reason (R) is false
  \item (D) Assertion (A) is false but Reason (R) is true
\end{enumerate}
\vspace{0.15cm}

\question[1] Assertion (A): To construct a pair of tangents to a circle from an external point, the angle between the tangents is supplementary to the angle subtended by the line segment joining the points of contact at the center. Reason (R): The sum of opposite angles in a cyclic quadrilateral is $180^\circ$.
\begin{enumerate}[label=(\Alph*)]
  \item (A) Both Assertion (A) and Reason (R) are true and Reason (R) is the correct explanation of Assertion (A)
  \item (B) Both Assertion (A) and Reason (R) are true but Reason (R) is NOT the correct explanation of Assertion (A)
  \item (C) Assertion (A) is true but Reason (R) is false
  \item (D) Assertion (A) is false but Reason (R) is true
\end{enumerate}
\vspace{0.15cm}

\question[1] Assertion (A): It is impossible to construct a tangent to a circle from a point lying inside the circle. Reason (R): A tangent touches the circle at only one point.
\begin{enumerate}[label=(\Alph*)]
  \item (A) Both Assertion (A) and Reason (R) are true and Reason (R) is the correct explanation of Assertion (A)
  \item (B) Both Assertion (A) and Reason (R) are true but Reason (R) is NOT the correct explanation of Assertion (A)
  \item (C) Assertion (A) is true but Reason (R) is false
  \item (D) Assertion (A) is false but Reason (R) is true
\end{enumerate}
\vspace{0.15cm}

\question[1] Assertion (A): The number of tangents that can be drawn from a point on the circle to the circle is exactly two. Reason (R): A point on the circle lies on the circumference.
\begin{enumerate}[label=(\Alph*)]
  \item (A) Both Assertion (A) and Reason (R) are true and Reason (R) is the correct explanation of Assertion (A)
  \item (B) Both Assertion (A) and Reason (R) are true but Reason (R) is NOT the correct explanation of Assertion (A)
  \item (C) Assertion (A) is true but Reason (R) is false
  \item (D) Assertion (A) is false but Reason (R) is true
\end{enumerate}
\vspace{0.15cm}

\question[1] Assertion (A): When constructing tangents to a circle from an external point $P$, we bisect the line segment $OP$ where $O$ is the center. Reason (R): The perpendicular bisector of $OP$ helps in finding the midpoint which acts as the center of a circle passing through $O$ and $P$.
\begin{enumerate}[label=(\Alph*)]
  \item (A) Both Assertion (A) and Reason (R) are true and Reason (R) is the correct explanation of Assertion (A)
  \item (B) Both Assertion (A) and Reason (R) are true but Reason (R) is NOT the correct explanation of Assertion (A)
  \item (C) Assertion (A) is true but Reason (R) is false
  \item (D) Assertion (A) is false but Reason (R) is true
\end{enumerate}
\vspace{0.15cm}

\end{questions}
\bigskip
\noindent\textbf{\large Section B: Very Short Answer (2 marks each)}

\begin{questions}
\question[2] If $\tan A = \frac{4}{3}$, find the value of $\sin A + \cos A$.
\vspace{0.15cm}

\question[2] Evaluate: $\frac{2 \tan^2 45^\circ + \cos^2 30^\circ - \sin^2 60^\circ}{1}$
\vspace{0.15cm}

\question[2] If $\sin(A - B) = \frac{1}{2}$ and $\cos(A + B) = \frac{1}{2}$, where $0^\circ < A + B \le 90^\circ$ and $A > B$, find $A$ and $B$.
\vspace{0.15cm}

\question[2] Express $\cos A$ in terms of $\tan A$.
\vspace{0.15cm}

\question[2] Evaluate $\frac{\tan 65^\circ}{\cot 25^\circ}$.
\vspace{0.15cm}

\question[2] Draw a circle of radius $3 \text{ cm}$. Take two points $P$ and $Q$ on one of its extended diameter each at a distance of $7 \text{ cm}$ from its centre. Draw tangents to the circle from these two points $P$ and $Q$.
\vspace{0.15cm}

\question[2] Construct a tangent to a circle of radius $4 \text{ cm}$ from a point on the concentric circle of radius $6 \text{ cm}$ and measure its length.
\vspace{0.15cm}

\question[2] Draw a line segment $AB$ of length $8 \text{ cm}$. Taking $A$ as centre, draw a circle of radius $4 \text{ cm}$ and taking $B$ as centre, draw another circle of radius $3 \text{ cm}$. Construct tangents to each circle from the centre of the other circle.
\vspace{0.15cm}

\question[2] Draw a circle of radius $3.5 \text{ cm}$. Construct a pair of tangents to this circle which are inclined to each other at an angle of $60^\circ$.
\vspace{0.15cm}

\question[2] Draw a circle of radius $5 \text{ cm}$. Mark a point $P$ at a distance of $8 \text{ cm}$ from the centre. Construct the two tangents from $P$ to the circle.
\vspace{0.15cm}

\end{questions}
\bigskip
\noindent\textbf{\large Section C: Short Answer (3 marks each)}

\begin{questions}
\question[3] If $\tan(A+B) = \sqrt{3}$ and $\tan(A-B) = \frac{1}{\sqrt{3}}$, where $0^\circ < A+B \le 90^\circ$ and $A > B$, find the values of $A$ and $B$.
\vspace{0.15cm}

\question[3] If $3\cot \theta = 2$, evaluate the value of $\frac{4\sin \theta - 3\cos \theta}{2\sin \theta + 6\cos \theta}$.
\vspace{0.15cm}

\question[3] Prove that $(\sin A + \csc A)^2 + (\cos A + \sec A)^2 = 7 + \tan^2 A + \cot^2 A$.
\vspace{0.15cm}

\question[3] If $\sin \theta + \sin^2 \theta = 1$, prove that $\cos^2 \theta + \cos^4 \theta = 1$.
\vspace{0.15cm}

\question[3] Evaluate: $\frac{\cos 45^\circ}{\sec 30^\circ + \csc 30^\circ}$.
\vspace{0.15cm}

\question[3] Draw a circle of radius $3 \text{ cm}$. Take two points $P$ and $Q$ on one of its extended diameters each at a distance of $7 \text{ cm}$ from its centre. Draw tangents to the circle from these two points $P$ and $Q$.
\vspace{0.15cm}

\question[3] Construct a tangent to a circle of radius $4 \text{ cm}$ from a point on the concentric circle of radius $6 \text{ cm}$ and measure its length.
\vspace{0.15cm}

\question[3] Draw a line segment $AB$ of length $8 \text{ cm}$. Taking $A$ as center, draw a circle of radius $4 \text{ cm}$ and taking $B$ as center, draw another circle of radius $3 \text{ cm}$. Construct tangents to each circle from the center of the other circle.
\vspace{0.15cm}

\question[3] Draw a circle of radius $3.5 \text{ cm}$. Construct a pair of tangents to this circle which are inclined to each other at an angle of $60^\circ$.
\vspace{0.15cm}

\question[3] Construct a triangle with sides $5 \text{ cm}, 6 \text{ cm}$ and $7 \text{ cm}$. Then construct a circle touching the side of $6 \text{ cm}$ at its midpoint and passing through the other vertex.
\vspace{0.15cm}

\end{questions}
\bigskip
\noindent\textbf{\large Section D: Long Answer (4-5 marks each)}

\begin{questions}
\question[5] Prove that $\frac{\cos A - \sin A + 1}{\cos A + \sin A - 1} = \csc A + \cot A$, using the identity $\csc^2 A = 1 + \cot^2 A$.
\vspace{0.15cm}

\question[5] If $\sin \theta + \cos \theta = \sqrt{3}$, then prove that $\tan \theta + \cot \theta = 1$.
\vspace{0.15cm}

\question[5] Evaluate: $\frac{5 \cos^2 60^\circ + 4 \sec^2 30^\circ - \tan^2 45^\circ}{\sin^2 30^\circ + \cos^2 30^\circ}$
\vspace{0.15cm}

\question[5] If $\sec \theta = x + \frac{1}{4x}$, prove that $\sec \theta + \tan \theta = 2x$ or $\frac{1}{2x}$.
\vspace{0.15cm}

\question[5] Draw a circle of radius $3 \text{ cm}$. Take two points $P$ and $Q$ on one of its extended diameters each at a distance of $7 \text{ cm}$ from its centre. Draw tangents to the circle from these two points $P$ and $Q$.
\vspace{0.15cm}

\question[5] Construct a triangle $ABC$ with side $BC = 6 \text{ cm}$, $AB = 5 \text{ cm}$ and $\angle ABC = 60^\circ$. Then construct a triangle whose sides are $\frac{3}{4}$ of the corresponding sides of the triangle $ABC$. Finally, construct a tangent to the circumcircle of triangle $ABC$ at vertex $A$.
\vspace{0.15cm}

\question[5] Draw a pair of tangents to a circle of radius $4 \text{ cm}$, which are inclined to each other at an angle of $60^\circ$. Justify the construction.
\vspace{0.15cm}

\question[5] Draw two concentric circles of radii $3 \text{ cm}$ and $5 \text{ cm}$. Construct a tangent to the smaller circle from a point on the larger circle and measure its length. Verify the measurement by calculation.
\vspace{0.15cm}

\end{questions}
\bigskip

\end{document}
